% ============================================================
% 第2章 数据概览与版本管理
% 【负责人：数据清洗与基础特征线（已按 2.0 版实测预填，可按需增删）】
% ============================================================
\section{数据概览与版本管理}

\subsection{输入数据概览}

官方提供四大类输入数据，声明时间范围为 2026-06-01 至 2026-07-30，数据实际
覆盖 2026-06-01 00:00:02 至 2026-07-31 23:59:59，规模约数十 GB：

\begin{table}[H]
\centering
\caption{输入数据一览（source\_version = 2.0 实测）}
\label{tab:input-data}
\small
\begin{tabularx}{\textwidth}{lLl}
\toprule
数据类别 & 内容与关键实测 & 文件形态 \\
\midrule
车辆画像 & 500 台车一行：能源类型、月均行驶里程/时长、月平均停留次数、
高速里程占比、早晨/黄昏/夜间行驶时长占比与夜间里程占比；无缺失；
能源构成为电力 411 / 柴油 80 / 天然气 9 & \texttt{.xlsx} \\
风险事件 & 24 类风险事件编码，原始 3{,}152{,}654 条、覆盖 500 台车；
含事件类型、事件时间、速度、经纬度；缺失集中在经纬度（4{,}933 条） &
单个 gz 分片，流式解析 \\
车辆轨迹 & 速度、里程、经纬度等轨迹流，30 个分片 & 分片文件 \\
车载 IMU & 加速度三轴 + EMS/GPS 双通道速度等信号，30 个分片 & 分片文件 \\
\bottomrule
\end{tabularx}
\end{table}

官方说明该数据集来自真实驾驶数据、但经过筛选、调整与模糊化处理，在比例和特征上
不能反映真实群体；因此本队所有建模结论均以窗口内自算统计量为准，不依赖画像
统计期不明带来的外部先验（见 4.1 节团队决策）。

\subsection{数据版本管理：从 1.0 到 2.0}

官方于 2026-09-17 发布 2.0 版数据，与早期 1.0 版相比发生实质性变化：画像名单
变化 157 台（31.4\%），事件行数减少 21.8\%（4{,}029{,}571 → 3{,}152{,}654），
正类车数由 149 变为 126（主口径）。团队于 2026-09-18 完成 2.0 版独立核对
（画像与事件设备号交集 500、双侧余量 0），并做出以下版本管理约定：

\begin{itemize}
  \item \textbf{全组统一切换 2.0 版}：目标名单即 2.0 版画像的 500 台车；
  \item 1.0 版全部产物标记 \texttt{source\_version=1.0} 归档留存、不删除，
        但\textbf{严禁}与 2.0 版数字混入同一实验或同一比较表；
  \item 四项版本号贯穿全部产物：\texttt{source\_version}（数据版本）、
        \texttt{cleaning\_version}（清洗版本）、\texttt{feature\_version}
        （特征版本）、\texttt{split\_version}（划分版本），任何实验偏离
        统一配置必须在实验登记中显式声明差异。
\end{itemize}

\begin{TODO}
2.0 版 IMU 与轨迹数据官方分享链接返回 403（InvalidAccessKeyId），截至本文撰写
尚未获取；第 4 章中 F1（轨迹暴露量）与 F2（IMU 特征）两块按“结构验证基于 1.0
分片、数值结论待 2.0 数据”的口径撰写。若提交前已拿到 2.0 传感器数据，请数据线
更新本节与第 4 章相应表述。
\end{TODO}

\subsection{时间窗口与关键口径}

团队于 2026-09-18 冻结的主口径为 \texttt{main\_20\_40}：
\begin{itemize}
  \item 有效数据窗口 $[\text{2026-06-01},\ \text{2026-07-31})$，左闭右开；
  \item 前 20 天（6/1--6/21）为特征窗口，后 40 天（6/21 起）为标签窗口；
  \item 2026-07-31 当天记录（59{,}357 条）因官方口径存疑整日\textbf{隔离保存}
        （不删除、不进特征与标签统计），待官方答复后可逆地并入或弃用；
  \item 正式推理窗：以 7/31 为截点，取最近 20 天同定义特征，对 500 台车
        全量输出预测。
\end{itemize}
30 天预测跨度仅作敏感性检查（\texttt{aux\_30d\_sensitivity}），不与主口径混用。
